% Generated draft: requires mathematical and source review.
\documentclass[11pt]{article}
\usepackage[margin=25mm]{geometry}
\usepackage{fontspec}
\setmainfont{Latin Modern Roman}
\usepackage{amsmath,amssymb,mathtools,unicode-math}
\setmathfont{Latin Modern Math}
\usepackage{enumitem,needspace,xurl,hyperref}
\setlength{\parindent}{0pt}
\setlength{\parskip}{5pt}
\setlength{\emergencystretch}{3em}
\newcommand{\N}{\mathbb{N}}
\newcommand{\Z}{\mathbb{Z}}
\newcommand{\Q}{\mathbb{Q}}
\newcommand{\R}{\mathbb{R}}
\newcommand{\C}{\mathbb{C}}
\newcommand{\F}{\mathbb{F}}
\newcommand{\A}{\mathbb{A}}
\newcommand{\PP}{\mathbb{P}}
\newcommand{\E}{\mathbb{E}}
\newcommand{\eps}{\varepsilon}
\newcommand{\dd}{\,\mathrm{d}}
\begin{document}
\section*{UnsolvedMath ID: 30004488}
\textbf{Source problem code:} catalogue entry
\subsection{Definitions and mathematical statement}
Let \(G\) be a finite group and let \(p\) be an odd prime. For a group homomorphism \(f\colon A\to B\), its kernel is \(\ker f=\{a\in A:f(a)=0_B\}\), using additive notation. The supplied record does not define \(C_{ba+}(G,p)\), the map \(\Psi_G\), its domain or codomain, or the group \(D_\Omega(G)\). In particular, it does not identify an ambient group in which the proposed equality is to hold.

\textbf{Question as supplied.} Is it true that, for every finite group \(G\) and every odd prime \(p\),
\[
C_{ba+}(G,p)=\ker\Psi_G\,?
\]
Does the assignment \(G\mapsto D_\Omega(G)\) admit a well-defined biset-functor structure?

A finite \((H,G)\)-biset is a finite set with commuting left \(H\)- and right \(G\)-actions. In the usual integral biset category, morphisms from \(G\) to \(H\) are formal integer combinations of isomorphism classes of such bisets, subject to the disjoint-union relations; composition is induced by the balanced product over the intermediate group. A biset functor on this category is an additive functor to abelian groups. This describes the standard terminology, but the supplied record does not specify whether this full category, a category restricted to particular groups or bisets, or another coefficient convention is intended. It also does not specify the proposed maps on \(D_\Omega\) or whether \(p\) is an implicit parameter.

\textbf{Unresolved formulation qualifications.} The displayed equality and the functor question retain the supplied targets, but cannot be made self-contained without the missing definitions and conventions. No restriction of \(G\) to a \(p\)-group is supplied. The background reports an established exact sequence for \(p=2\), but supplies neither that sequence nor its hypotheses; this is contextual source information, not an additional theorem asserted here.
\subsection{Short English statement}
For every finite group \(G\) and odd prime \(p\), the supplied problem asks whether \(C_{ba+}(G,p)\) equals the kernel of \(\Psi_G\), and whether the groups \(D_\Omega(G)\) form a well-defined biset functor. The objects, maps, and intended category are insufficiently specified in the record, so this is a statement-only draft requiring mathematical review.
\subsection{Sources}
\begin{itemize}
\item Supplied local problem registry: permanent UnsolvedMath ID \(30004488\), source problem code \texttt{catalogue entry}, title \emph{Biset-Functor Structure of Dade-Type Groups}, collection \emph{Oberwolfach Reports Open Problems}. No individual proposer is identified; the supplied proposed year is 2020.
\item Supplied report citation: \emph{Cohomology of Finite Groups: Interactions and Applications} (2021), Oberwolfach Reports, pp.\ 1197--1261, DOI \texttt{10.4171/owr/2020/23}. \url{https://doi.org/10.4171/owr/2020/23}. The report is described as documenting workshop discussions and open problems.
\item Supplied literature source: OWR Cohomology of Finite Groups report. \url{https://ems.press/content/serial-article-files/46862}.
\item Supplied literature source: Bouc--Yalcin, \emph{Borel--Smith functions and the Dade group}. \url{https://www.lamfa.u-picardie.fr/bouc/boya.pdf}. This attribution is supplied for the reference, not for the proposal of this problem.
\item Supplied literature assessment dated \texttt{2026-08-22}: the record labels the problem open, reports that the \(p=2\) exact sequence is established, and states that the odd-prime equality and corresponding biset-functor question remain open, with no later general resolution found in that review. This is supplied dated curation, not an independently verified literature assessment. The record itself requests independent verification of time-sensitive status and cited claims.
\end{itemize}
\end{document}
